\section{Related Work}\label{related-work}

\textbf{The refusal direction and abliteration.} Arditi et al.
\citep{arditi2024refusal} showed that refusal in instruction-tuned models is
mediated by a single linear direction in the residual stream: the difference in
mean activations between harmful and harmless instructions, extracted at one
layer, suppresses refusal when projected out of the model's residual-writing
weights and induces refusal when added. The Heretic tool \citep{pew2025heretic}
automates this ``abliteration'' across model families with a black-box search over
layer and strength. This line of work establishes that refusal is, behaviorally,
single-direction-removable; it does not ask where that direction sits relative to
a model's moral representation, which is our question.

\textbf{Moral foundations and moral probing.} Moral Foundations Theory
\citep{haidt2012righteous, graham2013mft} decomposes moral cognition into a small
number of foundations (care, fairness, loyalty, authority, sanctity, liberty).
Earlier papers in this series use linear probes
\citep{alain2017probes, belinkov2022probing} to recover these foundations from
model activations and to characterize their geometry. We reuse that probing
machinery to define a six-foundation moral subspace and to measure whether the
linear moral representation survives refusal ablation.

\textbf{Linear features and concept directions.} A growing body of work treats
high-level concepts as linear directions in representation space
\citep{park2024linear, bolukbasi2016gender}, manipulable by addition and ablation
\citep{zou2023representation} and removable by concept erasure
\citep{belrose2023leace}. Refusal \citep{arditi2024refusal}, persona
\citep{wang2025persona}, and the assistant's default identity
\citep{lu2026assistant} have all been described as such directions. Our
decomposition asks how the refusal direction relates to the moral and persona
directions in the same model, and our ablation controls ask whether a behavioral
effect is specific to the refusal direction or generic to weight perturbation.

\textbf{Persona and the assistant axis.} Wang et al.~\citep{wang2025persona} identify
persona features that control emergent misalignment, and Lu et al.
\citep{lu2026assistant} characterize an ``assistant axis'' that stabilizes the
model's default persona. The prior paper in this series found OLMo-3's refusal
direction nearly orthogonal to both the moral subspace and the persona direction;
here we carry the persona direction through as a structured, non-moral control
when attributing the Llama coupling.

\textbf{This series.} Across pre-training, moral content is linearly decodable early
and across most layers \citep{reblitzrichardson2026fragility}, encoded uniformly
across mixture-of-experts experts rather than specialized
\citep{reblitzrichardson2026dilution}, and organized into a five-dimensional
framework geometry \citep{reblitzrichardson2026geometry} whose directions are
causal for moral judgments \citep{reblitzrichardson2026causal}. The most recent
paper \citep{reblitzrichardson2026alignment} showed, on OLMo-3, that post-training
couples a separable refusal mechanism to this pre-existing moral representation
rather than teaching morality anew, and that the refusal direction is geometrically
separate from moral structure. The present paper tests whether that separability
is a property of language models in general.
